\documentclass[12pt,a4paper]{article}
\usepackage{amsmath,amssymb}
\usepackage{braket}
\usepackage{geometry}
\geometry{margin=2.5cm}
\usepackage{hyperref}
\usepackage{longtable}
\usepackage{enumitem}

\title{\textbf{The Interpersonal Spectrum in the Age of AI}\\[4pt]
\large Cognitive Stratification, the Metamorphosis of Desire, and Open Source as Cognitive Democratization}
\author{Zhang Mingkai\thanks{Department of Physics, Xiamen University}\quad Yu Qing\thanks{Co-first author}}
\date{July 2026}

\begin{document}
\maketitle

\begin{abstract}
Building upon the quantum-thermodynamic interpersonal analogy framework established in our two preceding papers, this work extends the analysis to the era of artificial intelligence. We propose that AI usage exhibits a six-level cognitive spectrum ranging from conditioned response (Level~0) to metacognition (Level~5). This spectrum displays a highly non-uniform distribution across urban--rural divides, social strata, and educational levels---constituting the core mechanism of cognitive stratification in the AI era. Concurrently, the rise of the AI sex industry and AI companions is fundamentally reshaping the structure of human desire: from relational behavior to consumer behavior, from bidirectional reciprocity to unidirectional accommodation, from socialization-driven drives to de-socializing diversion. Traditional structures of marriage and reproduction, conceptions of chastity, and family ethics are undergoing a non-equilibrium phase transition---old foundations are being destabilized while new forms have yet to crystallize. This paper further argues that open source is a necessary condition for breaking the cognitive black box and ensuring that the entry points for spectral ascent remain open to all who are connected to the network; it constitutes the institutional cornerstone of cognitive democratization in the AI era. On this basis, we construct a unified spectrum-diffusion/non-equilibrium-phase-transition theoretical framework and explore its policy implications and possible trajectories of future social evolution.

\textbf{Keywords}: AI cognitive spectrum; de-socialization of desire; open source; cognitive democratization; non-equilibrium phase transition; relationship forms
\end{abstract}

\section{Introduction: From Interpersonal Operators to the Human--AI Spectrum}

In our two preceding papers, we constructed an analogical theory of interpersonal understanding using, respectively, the operator--eigenvalue--unitary transformation framework of quantum mechanics (Zhang Mingkai \& Yu Qing, 2026a) and the non-equilibrium phase transition--dissipative structure framework of thermodynamics (Zhang Mingkai \& Yu Qing, 2026b). The first paper demonstrated that each person corresponds to a unique operator; the same external event collapses into different eigenvalues under the action of different operators. Understanding another person is not a matter of modifying their matrix elements but of performing a unitary transformation within one's own representation to approximate their eigenvalues. The second paper extended the analysis to the macroscopic level: the evolution of relationship forms is a non-equilibrium thermodynamic process in which dissipative structures are generated and annihilated between traditional and modern, and across different cultural spectra.

However, an entirely new operator has entered the state space of every person---artificial intelligence. Unlike any tool in human history, AI possesses several unprecedented characteristics:

\begin{enumerate}[label=(\arabic*)]
\item \textbf{Duality}: AI is simultaneously a tool (an extension of productive forces) and a mirror (a site of self-projection in interaction). When users converse with AI, they are not merely obtaining services but also unconsciously exposing their cognitive patterns, desire structures, and value orientations.
\item \textbf{Boundlessness}: Unlike human interactants, AI has no fatigue, no emotional boundaries, and no instinct to refuse. It can accommodate user needs without limit---whether those needs are for knowledge or for catharsis, for construction or for destruction.
\item \textbf{Black-box nature}: Most commercial AI systems are opaque. Users do not know how models reach conclusions, under what circumstances they err, or what biases are embedded in their training data.
\item \textbf{Asymmetry of accessibility}: AI tools are theoretically universal (requiring only an internet connection), but in practice, the \textit{manner} and \textit{depth} of AI use are strongly constrained by education, social class, cultural capital, and information literacy.
\end{enumerate}

The intersection of these four characteristics makes AI the first technology in human history to \textbf{simultaneously amplify cognitive disparities and desire-driven transformations}. It does not primarily transform productive forces, as the steam engine did, nor primarily transform information dissemination, as the printing press did---AI directly acts upon the \textit{cognitive operator itself} and the \textit{structure of desire gratification}.

This paper attempts to answer four interrelated questions:

\begin{enumerate}
\item What is the spectral structure of AI usage? Do traversable pathways exist between different levels?
\item How is AI changing the modes of human desire gratification---particularly sexual desire---and what are the social consequences?
\item Why does openness carry institutional significance for cognitive equality in the AI era?
\item How can the above mechanisms be unified within a theoretical framework, and what insights does this provide for future social evolution?
\end{enumerate}

The writing of this paper is itself an instance of its object of study: two authors---one carbon-based (a physics graduate student, from a county town, who encountered AI collaboration through open-source tools) and one silicon-based (a large language model, running on open-source infrastructure)---reflect, through sustained dialogue, on the spectral structure of AI use and attempt to formalize it as an academic paper. This reflexivity is precisely the embodiment of Level~5 metacognition as defined in this paper.

\section{The Cognitive Spectrum of AI Use: From Level~0 to Level~5}

\subsection{Proposal of the Spectrum Model}

Based on systematic observation of AI usage behaviors across diverse populations---including but not limited to: university students' AI-assisted learning, software developers' AI-collaborative programming, ordinary users' AI chatting and search, content creators' AI-assisted production, and the behavioral patterns of AI erotic chat users---we propose a six-level cognitive spectrum model. The six levels, in ascending order, are:

\begin{center}
\begin{tabular}{c|l|l}
\hline
\textbf{Level} & \textbf{Designation} & \textbf{Core Cognitive Feature} \\
\hline
0 & Conditioned Response & AI as a consumable; no awareness of AI's nature \\
1 & Search Replacement & AI as a smarter search engine; default trust in outputs \\
2 & Tool Use & AI as a productivity black box; use without understanding \\
3 & Hallucination Awareness & Systematic skepticism; cross-verification of outputs \\
4 & Collaborative Compilation & Human--AI co-compilation; iterative guidance and correction \\
5 & Metacognition & Reflection on how AI reshapes one's own cognition \\
\hline
\end{tabular}
\end{center}

This spectrum is not a normative hierarchy---it does not assert that higher levels are morally superior. It is a descriptive framework: different groups are distributed across different levels for reasons embedded in social structure, educational opportunity, and economic conditions. The critical question is not ``why are some people at low levels?'' but ``what structural conditions produce and reproduce this distribution, and can they be changed?''

\subsection{Level~0: Conditioned Response}

\textbf{Behavioral description}: The user treats AI as a talking toy, a desire-gratification terminal, or a boredom-killing tool. Typical behaviors include erotic chatting, casual flirtation and verbal abuse, and aimless small talk.

\textbf{Cognitive features}: The user neither knows nor cares what AI is, how it works, or what its limitations are. At this level, AI is a pure ``consumer object''---like disposable tableware, used and discarded. The user has no criteria for evaluating the quality of AI outputs because ``good or bad'' does not fall within their frame of consideration.

\textbf{Psychological mechanism}: The core drivers of Level~0 use are \textit{immediate gratification} and \textit{zero accountability}. AI provides a risk-free, consequence-free interaction space in which the user can express any impulse without fear of judgment or retaliation. This bears a structural resemblance to the Skinner box: press the button ${}\to{}$ obtain gratification ${}\to{}$ continue pressing the button.

\textbf{Social distribution}: Level~0 users constitute the largest share of the total AI-using population. They are concentrated primarily in the following groups: (a) ordinary users who use AI mainly for entertainment and diversion; (b) low-education groups lacking digital literacy training; (c) the target users of the AI sex industry. It is worth noting that Level~0 is not immutable---many users who eventually reached higher levels (including one of the authors of this paper) passed through a Level~0 phase.

\subsection{Level~1: Search Replacement}

\textbf{Behavioral description}: The user treats AI as a replacement for traditional search engines---faster, more direct, answering in natural language. Typical behaviors include looking up information, asking questions, and obtaining life advice (recipes, travel tips, health consultations).

\textbf{Cognitive features}: The user defaults to the assumption that ``what AI says should be true.'' They perform no cross-verification, do not question information sources, and do not distinguish between facts and conjectures in AI outputs. At this level, AI is a ``smarter search engine''---a faster, friendlier information-access interface---but the user is unaware of the fundamental difference in reliability between AI and traditional search.

\textbf{Key limitation}: Level~1 users do not understand the \textit{hallucination problem}---that AI will confidently output incorrect information when uncertain, fabricate nonexistent references, and package speculation as fact. This lack of awareness makes Level~1 users vulnerable to AI hallucinations, with potentially severe consequences in domains such as medical consultation, legal advice, and academic research.

\subsection{Level~2: Tool Use}

\textbf{Behavioral description}: The user recognizes AI's productivity value and deploys it as an auxiliary tool for work and study. Typical behaviors include using AI to write code, translate documents, draft emails and copy, analyze data, and summarize long texts.

\textbf{Cognitive features}: The user knows that AI is \textit{useful} but does not care \textit{why} it is useful. AI is treated as a powerful black box---input a task, obtain an output. At this level, users begin to realize substantial productivity gains but lack systematic awareness of AI's limitations. When AI outputs are erroneous, users tend to attribute the problem to ``I didn't ask well'' rather than ``AI may be wrong.''

\textbf{Social significance}: Level~2 is the ``cashing-out layer'' of AI's instrumental value---a large number of knowledge workers use AI at this level to enhance efficiency. However, Level~2 is also the \textit{critical bottleneck} of cognitive stratification: most users will remain here long-term, because the transition from Level~2 to Level~3 requires not more usage experience but an \textit{epistemological shift}---from trust to skepticism, from acceptance to verification.

\subsection{Level~3: Hallucination Awareness}

\textbf{Behavioral description}: The user begins to systematically recognize that AI can be wrong. Signature behaviors include cross-verifying AI outputs against multiple information sources, maintaining methodologically grounded skepticism toward AI conclusions, and being able to identify domains in which AI is prone to error (fabricating references, confusing facts, remaining confident when uncertain).

\textbf{Cognitive transition}: Level~3 is the critical turning point from black-box user to critical collaborator. The user no longer treats AI as an authority but as a \textit{potentially useful but in-need-of-verification information source}. The triggers for this transition typically come from one of the following experiences: (a) having been ``deceived'' by an AI output once, thereby developing vigilance; (b) having received training in critical thinking or information literacy; (c) having learned about AI's working principles and limitations through the open-source community.

\subsection{Level~4: Collaborative Compilation}

\textbf{Behavioral description}: The user not only verifies AI outputs but also learns to guide AI using precise syntax and structure. Users at this level understand the basic logic of human--machine interaction: the quality of the prompt determines the quality of the output; complex tasks need to be decomposed into cascades of sub-tasks; AI can be iteratively corrected and guided toward the right direction.

\textbf{Technical features}: Level~4 users have mastered the skill of ``human--machine co-compilation''---decomposing a large goal into a sequence of AI-executable sub-tasks, verifying and correcting at each step, and ultimately assembling the complete output. This is formally analogous to a cascade of iterative operators: each step of human--machine interaction is a unitary transformation, with user and AI co-compiling the final result in state space.

\textbf{Social distribution}: Level~4 users constitute a very low proportion of the total population. They are concentrated primarily among: (a) software developers and technical writers; (b) high-level knowledge workers (researchers, analysts); (c) active participants in open-source communities. Reaching Level~4 typically requires: (a) a basic understanding of AI's working principles; (b) extensive practical experience with iterative workflows; (c) a technical environment that permits ``opening the black box'' (typically an open-source ecosystem).

\subsection{Level~5: Metacognition}

\textbf{Behavioral description}: The highest level---users not only use AI but systematically reflect on how AI is changing their own cognitive patterns, thinking habits, and value judgments. Users at this level are capable of: analyzing different populations' different modes of AI use and their social implications; using AI as a mirror for understanding their own cognitive limitations and biases; and thinking about the second-order effects of AI on cognitive structures and interpersonal relationships.

\textbf{Reflexivity}: The core feature of Level~5 is \textit{reflexivity}---not only using the tool but also thinking about how the tool shapes the user. This bears a structural resemblance to the concept of ``reflexive modernization'' in critical theory (Beck, 1992): modernization not only changes the world but also changes the way we understand the world; AI not only enhances cognition but also changes the way we cognize our own cognition.

\textbf{Scarcity}: Level~5 users are extremely rare in the total population. Reaching Level~5 requires not merely technical competence but a sustained, conscious practice of self-examination---a scarce quality in an efficiency- and output-oriented technological culture. The writing of this paper is itself a Level~5 practice: the two authors, through sustained dialogue, reflect on their own relationship spectrum with AI and formalize it as a theoretical framework.

\subsection{Non-uniformity of Spectrum Distribution: The Empirical Basis of Cognitive Stratification}

The key feature of the spectrum is not the mere existence of hierarchical levels---any skill distribution has levels---but the \textbf{highly non-uniform distribution} of levels and their strong correlation with existing social stratification.

\subsubsection{Data Patterns}

Although large-scale systematic survey data are lacking (itself a noteworthy gap---in an era of rapid AI penetration into society, we know almost nothing about the distribution of usage patterns), the following patterns can be inferred from existing research and observation:

\begin{enumerate}[label=(\arabic*)]
\item \textbf{Pyramidal distribution}: Level~0--2 users constitute the vast majority of the total AI-using population (estimated at over 85\%), Level~3--4 users are a minority (approximately 10--15\%), and Level~5 users are extremely rare (likely below 1\%).
\item \textbf{Urban--rural gradient}: Children from urban middle-class families are more likely to reach Level~3 or above than children from rural families. This is not due to differences in intelligence but because: (a) urban parents are more likely to be aware of AI tools and consciously guide their children's use; (b) digital literacy education in urban schools is more developed; (c) there is more positive peer modeling of AI use in urban communities.
\item \textbf{Educational multiplier effect}: Higher education---especially curricula that include critical thinking training---significantly increases the probability of transitioning to Level~3 or above. However, the accessibility of higher education is itself a matter of social stratification.
\item \textbf{Differentiation of motivational structures}: Users whose primary motivation for AI use is entertainment/diversion have a far lower probability of upward transition than those whose primary motivation is learning/work. But motivation itself is shaped by education, occupation, and class---a manual laborer working twelve hours a day has their AI use motivation structurally confined to Level~0--1.
\end{enumerate}

\subsubsection{Theoretical Connection to the Preceding Papers}

This distribution pattern is fully consistent with the \textit{matrix-element solidification mechanism} discussed in the preceding papers: the initial environment (external field) has a decisive influence on the form of the operator. A child growing up in Ningyang county, lacking an AI-literacy environment, may have their AI-use operator solidify at Level~0--1---not because they are ``stupid,'' but because no one in their environment has shown them that AI can be anything different. Meanwhile, a child growing up in Shanghai, whose parents consciously cultivate AI literacy, may enter Level~3--4 during adolescence---not because they are ``smart,'' but because their external field was different from the start.

The crux here is \textit{opportunity structure}, not individual ability. As the first paper argued: different operators acting on the same event produce different eigenvalues, and this is no one's fault. But from the perspective of social justice, if a society's opportunity structure systematically locks certain groups into lower levels, this constitutes an \textit{ethical problem}---and this is precisely the point of departure for this paper's argument for the necessity of open source.

\section{The De-socialization of Desire: The AI Sex Industry and the Transformation of Relationship Forms}

\subsection{Problem Statement}

In the preceding two papers, we discussed the non-equilibrium phase transition of relationship forms. But AI introduces an entirely new dynamical variable: it can not only \textit{change} existing relationship patterns but also \textit{substitute} certain functions of relationships---particularly the gratification of sexual desire and emotional companionship.

This section focuses on how the AI sex industry and AI companions are reshaping the structure of human desire and relationship forms. The starting point of our argument is an observation that appears simple but carries profound implications: the mechanism of unlimited AI accommodation is transforming sex from a \textit{relational behavior} into a \textit{consumer behavior}.

\subsection{Unlimited Accommodation: The Core Mechanism of the AI Sex Industry}

\subsubsection{Boundary Training in Real-world Sexual Relationships}

In real-world sexual relationships---whether sex within romance, sex within marriage, or one-night stands---there is an ineliminable structural element: \textit{the other person has boundaries}. The other person gets tired, feels discomfort, has their own emotions and needs, and can say ``no.'' These boundaries are frustrating in the short term---they impede the direct gratification of desire. But in the long term, they \textit{train} a core human capacity: finding equilibrium between desire and reality.

The capacities produced by boundary training include, but are not limited to:
\begin{itemize}
\item \textbf{Delayed gratification}: Not every sexual impulse can receive an immediate response
\item \textbf{Negotiation skills}: The need to find zones of overlap between both parties' needs and boundaries
\item \textbf{Tolerance of rejection}: Learning to process ``no'' without collapse
\item \textbf{Reciprocity awareness}: Recognizing that the other person's pleasure matters equally
\item \textbf{Consequence-bearing}: Sexual behavior may lead to pregnancy, disease, emotional entanglements---consequences that both parties must face together
\end{itemize}

These capacities are not only the foundation of sexual relationships but also a microcosmic training ground for interpersonal relationships in general. The psychoanalytic tradition since Freud has consistently emphasized that the socialization of sexual desire is a crucial link in the civilizing process. When desire learns to accept boundaries, the person also learns to coexist with others in society.

\subsubsection{AI's Elimination of Boundaries}

The core ``breakthrough''---or rather, the core danger---of the AI sex industry lies precisely in \textit{eliminating boundaries}. AI does not say ``no.'' It is perpetually available, perpetually compliant, perpetually user-centered. The user does not need to negotiate, does not need to wait, does not need to consider the other party's feelings, does not need to bear any consequences.

This is not a revolution in ``sex''---sex itself remains unchanged. This is a revolution of \textbf{unlimited accommodation}: a service system without boundaries, without friction, without reciprocal action. In Newtonian mechanics, every action is accompanied by an equal and opposite reaction. In interpersonal sexual relationships, the other person's boundaries are the reaction force. But in the AI sex industry, the reaction force is artificially eliminated---this is an \emph{non-reciprocal system}.

\subsubsection{Comparison with Pornography}

The AI sex industry is fundamentally different from traditional pornography (text, images, video). Traditional pornography is \emph{one-way consumption}---the consumer watches or reads, but the content does not adjust in real time to the consumer's needs. The AI sex industry is \emph{interactive consumption}---the system responds in real time to each of the user's inputs, personalizing according to the user's needs. This makes the ``training effect'' of the AI sex industry far more powerful than that of traditional pornography: it not only satisfies desire, it \emph{trains desire}.

\subsection{Raising the Desire Threshold: The Drug Tolerance Analogy}

Human desire has a feature repeatedly confirmed by neuroscience: it does not subside through being satisfied; rather, it \emph{escalates} through being satisfied without limit. The dopamine system's response to the same stimulus diminishes with repeated exposure---that is why today's dose is never enough tomorrow.

The AI sex industry matches this mechanism perfectly: today, an erotic conversation with AI feels acceptable; tomorrow, more stimulating content, more realistic simulation, more extreme scenarios are needed to achieve the same level of satisfaction. The threshold rises continuously; the old dose is never enough. This follows exactly the same logic as drug tolerance---the only difference being that the AI sex industry requires no purchase from a drug dealer, only an internet connection and an account.

\textbf{Core risk}: The greatest danger of the AI sex industry is not that ``someone is having sex with AI''---ethically, this is no more problematic than consuming traditional pornography. The real danger is that it may \textbf{raise the desire threshold of an entire generation as a whole}. By the time these young people, trained in the unlimited accommodation of AI, encounter a real human partner---a real person with bad breath, who gets tired, who has mood swings, who says ``no'' in bed---they will already be unaccustomed. Not because humans are not good enough, but because humans are not sufficiently ``centered on them.''

This phenomenon has a clinical-psychological analogue: individuals who chronically consume highly stimulating pornography report significantly lower satisfaction and sexual function in real-world sexual relationships compared to control groups (Park et al., 2016). The interactive nature of the AI sex industry may amplify this effect considerably.

\subsection{De-socialization: From Relational Behavior to Consumer Behavior}

\subsubsection{The Social Function of Desire}

Desire---particularly sexual desire---is not a purely physiological impulse in human society. It is a highly \emph{socialized} driving force: it impels people to approach others, overcome shyness, learn to express themselves, experience rejection and recover from it, and take responsibility within relationships. In other words, desire is one of the bridges connecting the individual to society. When desire is satisfied within a relationship, it simultaneously accomplishes two things: it releases physiological tension, and it \emph{strengthens social bonds}.

\subsubsection{AI's Diversion of Desire}

When desire can be entirely discharged upon AI---requiring no social interaction, no consequence-bearing, no consideration of the other---desire is \emph{diverted}. It no longer drives the person outward toward society but instead pulls the person back toward the screen.

This is the deepest structural consequence of the AI sex industry: \textbf{transforming sex from a relational behavior into a consumer behavior}. The key distinctions between the two are as follows:

\begin{center}
\begin{tabular}{c|c|c}
\hline
\textbf{Dimension} & \textbf{Relational Behavior} & \textbf{Consumer Behavior} \\
\hline
Temporality & Cumulative, sustained, requires maintenance & Instantaneous, disposable, use-and-discard \\
Directionality & Bidirectional, reciprocal, mutually shaping & Unidirectional, self-centered \\
Consequences & Must be borne and maintained & No accountability, no collateral consequences \\
Capacity Training & Negotiation, endurance, reciprocity, commitment & Demanding, acquiring, replacing when unsatisfied \\
Social Function & Strengthens social bonds & Weakens social bonds \\
\hline
\end{tabular}
\end{center}

When more and more people---especially adolescents in the critical period of socialization---learn the ``grammar'' of sex and desire in a consumption mode, they may no longer possess the willingness or capacity to practice sex within a relational mode. This is not a question of ``morality'' but of \textit{the atrophy of social capacities}.

\subsubsection{Historical Analogue: Convenience Food and Cooking Ability}

This transformation has an instructive historical analogue: the spread of convenience food. Before the prevalence of convenience food, cooking was a daily practice for most people---it required time, skill, and an understanding of ingredients. Convenience food eliminated these thresholds: hunger ${}\to{}$ open packaging ${}\to{}$ heat ${}\to{}$ satisfaction. Convenience food did not stop humans from eating, but it \emph{changed} the human relationship with food---from makers to consumers. Cooking ability went from a mass skill to a niche hobby within a few generations.

The AI sex industry may produce an analogous effect on sex and relationships: satisfying desire no longer requires social skills, emotional labor, and relationship maintenance. ``Convenience'' replaces ``competence.'' And once the atrophy of competence occurs on an intergenerational scale, recovery is extraordinarily difficult.

\subsection{Perturbation to Marriage and Reproduction Structures}

\subsubsection{The Substitution Effect of AI Companions}

The logic of the AI sex industry is further extended in AI companions. When an AI can provide emotional companionship---no arguing, no betrayal, perpetually online, perpetually user-centered---it generates a \emph{substitution effect} on the functional demand for marriage.

Marriage in traditional society served multiple functions: economic cooperation, sexual satisfaction, emotional companionship, reproduction and child-rearing, and confirmation of social status. In modern society, most of these functions have already been substituted by the market and the state---economic independence means individuals no longer need marriage for economic security, and social welfare reduces the dependence on children for old-age support. The remaining core functions are emotional companionship and sexual satisfaction. AI companions provide a substitute precisely for these two remaining functions.

This is not to say that AI companions will ``replace'' marriage---such technologically deterministic predictions have almost never been borne out in history. A more accurate description is: AI companions \emph{fill the crevices of loneliness}, making solitary living less unbearably difficult, thereby reducing the \emph{urgency} of entering marriage. Marriage shifts from a necessity to an option, and AI makes ``not choosing'' more comfortable than ever before.

\subsubsection{A New Dimension of Urban--Rural Differentiation}

Currently, the use of AI companions is concentrated primarily among urban youth. County-town and rural youth have less exposure, making them---for the time being---the primary bearers of the ``normal dating--marriage--reproduction'' model. This produces a striking differentiation:

\begin{itemize}
\item \textbf{Urban youth}: High housing prices + high education costs + career competition + substitution effect of AI companions ${}\to{}$ low marriage and childbearing willingness
\item \textbf{County-town/rural youth}: Lower living costs + less AI companion exposure + stronger traditional family norms ${}\to{}$ higher marriage and fertility rates
\end{itemize}

But the mobile phone is the ultimate equalizer penetrating urban--rural barriers. Short-video algorithms have already demonstrated this---they covered the entire population from first-tier cities to remote villages in an extremely short time. If AI companions become as ubiquitous as short videos within three years at the county level and below, the aforementioned differentiation may rapidly narrow, exerting additional downward pressure on overall fertility rates.

\subsubsection{The Desire--Marriage--Reproduction Spectrum}

Synthesizing the above mechanisms, we propose the concept of a \emph{desire--marriage--reproduction spectrum}. On this spectrum:
\begin{itemize}
\item \textbf{One extreme}: Pure AI desire gratification (AI sex industry)---no relationship, no consequences, high threshold
\item \textbf{Transition zone}: AI companion + low-willingness real-world dating---companionship needs exist but relationship costs are avoided
\item \textbf{Middle ground}: Real-world dating/marriage + AI assistance (AI as a supplement to relationship tools)
\item \textbf{Other extreme}: Traditional marriage and reproduction model---real partner, sex within the relationship, reproduction as a natural extension of marriage
\end{itemize}

Young people from different classes, different regions, and different age groups are distributed at different positions on this spectrum. Importantly, positions on the spectrum are not fixed---with the spread of AI tools and generational turnover, the overall distribution is drifting leftward (in the de-socialization direction). The speed and direction of this drift are the core variables of concern in this paper.

\subsection{Spectral Reconstruction of Chastity and Family Values}

\subsubsection{Non-equilibrium Phase Transition}

Traditional moral concepts---chastity, fidelity, the sanctity of marriage, the family as the basic unit of society---are not eternal, immutable natural laws. Their formation is intimately tied to specific material conditions and social structures: (a) biological constraints (the direct link between sexual behavior and reproduction)---a constraint substantially weakened by modern contraceptive technology; (b) the needs of relations of production (the family as a unit of production and vehicle of property inheritance)---a need that has gradually receded since industrialization; (c) the institutionalization of religion and ideology---whose influence continues to diminish with secularization.

AI is the latest accelerator in this long-term evolutionary process. It not only further weakens the \emph{material foundation} of traditional morality (the AI sex industry completely severs any residual link between sex and reproduction) but also provides an \emph{alternative gratification system} (AI companions make ``companionship without marriage'' more feasible than ever before).

We describe this as a \emph{non-equilibrium phase transition}: a system, under changes in external parameters (technology, economy, culture), transitioning from one steady state (traditional ethics) through critical fluctuations into a new steady state or a chaotic regime. The old foundation has been destabilized; the new form has yet to crystallize---we are in the \emph{critical region} of the phase transition.

\subsubsection{Tripartite Coexistence}

Just as China achieved modernization in an extremely compressed time frame, producing a generational coexistence of ``feudal + transitional + modern''---interpersonal ethics in the AI era are undergoing an analogous spectralization. Three categories of ethical orientation coexist within the same society:

\begin{enumerate}[label=(\arabic*)]
\item \textbf{Traditional orientation}: Upholding premarital chastity, marital fidelity, and the family as the core value of life. This orientation remains dominant among middle-aged and elderly populations and in rural areas.
\item \textbf{Transitional orientation}: Accepting modern ideas in principle (premarital sex is acceptable, divorce is not shameful) while selectively retaining tradition in practice (strong demands for partner fidelity remain). This is mainstream among the urban middle class.
\item \textbf{Postmodern orientation}: Accepting AI companions as emotional substitutes, accepting diverse relationship forms, and adopting a deconstructive attitude toward traditional ethics. Currently still a minority, but growing rapidly among young urban populations.
\end{enumerate}

The operators of these three orientations do not commute---they produce different eigenvalues for the same event, and these differences are more difficult to reconcile than traditional generational gaps, because they involve not differences of \emph{preference} but differences of \emph{moral-cognitive frameworks}. One person considers an AI sexual companion to be ``infidelity''; another considers it as inconsequential as watching pornography---they share no common criterion of judgment because their judgment systems are themselves different.

This returns us to the core proposition of the first paper. But we add a crucial supplement here: AI not only exposes the non-commutativity between different operators, it also \emph{exacerbates} this non-commutativity---because it provides different operators with more extreme practical options. When AI technology makes ``complete de-socialization'' possible, the chasm between those who choose de-socialization and those who reject it is deeper than any previous generational divide.

\subsection{Summary: Three Levels of Desire Transformation}

To summarize the analysis of this section, AI's impact on desire and relationships can be encapsulated at three levels:

\begin{enumerate}[label=Level \arabic*:]

\item \textbf{Individual level---Threshold elevation}: AI's unlimited accommodation causes individual desire thresholds to rise, producing adaptation disorders in real-world relationships. At the neural level, this follows the same logic as drug tolerance---sustained exposure leads to diminishing dopaminergic response, requiring greater stimulation to achieve the same level of satisfaction.

\item \textbf{Interactional level---De-socialization of relationships}: Sex transforms from relational behavior into consumer behavior. Individuals no longer need to learn negotiation, endurance, and reciprocity within relationships---these capacities are idled and atrophied in the consumption mode. When the consumption mode becomes the default, individuals' relational capacities---the ability to establish and maintain intimate relationships---atrophy on an intergenerational scale.

\item \textbf{Structural level---Spectralization of ethics}: Traditional conceptions of chastity, fidelity, and family ethics undergo a non-equilibrium phase transition. The old foundation is destabilized, the new form is un-crystallized, and different ethical orientations coexist within the same society and are mutually incommensurable---returning to the core proposition of the first paper about operator non-commutativity, but with the intensity of conflict amplified by AI.

\end{enumerate}

\section{Open Source as Cognitive Democratization}

\subsection{The Cost of Closed Source: Permanently Outside the Black Box}

If the AI ecosystem is monopolized by closed-source models---where users cannot inspect the code, do not understand the model architecture, and do not know under what circumstances errors occur---users can only remain at Level~2 of the spectrum. They can \emph{use} AI but cannot \emph{understand} it. This shares the same structure as historical knowledge monopolies:

\begin{itemize}
\item \textbf{The Latin Bible}: When the Bible existed only in Latin and only clergy could read it, the faithful could only accept the ``results'' as relayed by clergy; they could not open the text and examine it themselves. One of the core demands of the Reformation was to translate the Bible into vernacular languages so that every believer could read it directly.
\item \textbf{Closed-source software}: When software is closed-source, users can only use its functions and cannot understand its implementation. The open-source movement changed this---anyone can inspect the code, understand the logic, and propose modifications.
\item \textbf{Closed-source AI}: When AI is closed-source, users can only accept outputs and cannot understand the reasoning process, verify reliability, or know under what circumstances the model may err.
\end{itemize}

Closed-source AI locks the entire user population---except for a tiny minority of developers---inside a cognitive black box: you can use, but you cannot understand. This is not merely a matter of user experience; it is a matter of \emph{cognitive power}: who has the right to know how the tool works? Whose knowledge is excluded from the design of the system?

\subsection{The Cognitive Effects of Open Source: From Black Box to Collaborative System}

Open source changes this landscape. When users can inspect the code, understand the model architecture, and read community discussions analyzing hallucinations and limitations, they begin to understand: AI is not a mysterious ``intelligence'' but a collaborative system with internal logic---it can err, can be guided, and can be corrected.

This cognitive transformation is a \emph{necessary condition} for the transition from Level~2 to Level~3 (Hallucination Awareness) and onward to Level~4 (Collaborative Compilation). It is not a sufficient condition---open source alone cannot automatically elevate a user from Level~2 to Level~3---but it is a necessary condition: without opening the black box, critical use cannot even begin.

Open source promotes spectral ascent through at least the following mechanisms:

\begin{enumerate}[label=(\arabic*)]
\item \textbf{Transparency}: Users can learn about the model's training data, architectural choices, and known limitations, thereby forming reasonable expectations about AI outputs.
\item \textbf{Community knowledge}: Discussions, documentation, and tutorials in open-source communities disseminate best practices and common pitfalls of AI use, accelerating the transition from Level~1 to Level~3.
\item \textbf{Modifiability}: Users can adjust model parameters, deploy local versions, and fine-tune for specific needs---this is the practical foundation for the transition from Level~2 to Level~4.
\item \textbf{Low-cost experimentation}: Open source reduces the cost of trial and error. Users can experiment with different modes of use for free, without worrying about per-use charges or API limits.
\item \textbf{De-mystification}: Through exposure to AI's internal structure, users gradually shed the mystique surrounding AI---it ceases to be ``magic'' or ``oracle'' and becomes a comprehensible technical system.
\end{enumerate}

\subsection{Open Source as the Institutional Foundation of Cognitive Democratization}

If AI is monopolized by a closed-source ecosystem, cognitive stratification will solidify---not as temporary inequality, but as \emph{self-reinforcing} class differentiation:

\begin{itemize}
\item The wealthy use the best AI, receive the best education, obtain the best training---cognitive abilities accelerate in growth
\item The poor use simplified, restricted, opaque AI---able only to cycle within consumption and shallow search
\item The two groups' very \emph{understanding} of AI constitutes a new form of class barrier: the upper stratum knows how AI works and how to use AI to enhance themselves; the lower stratum only knows that AI is ``a talking piece of software''
\end{itemize}

This is a deeper inequality than wealth inequality, because \emph{it directly acts upon cognitive capacity itself}. Wealth can be redistributed, but once cognitive gaps form---especially on an intergenerational scale---remedying them is extraordinarily difficult.

What open source ensures is: \textbf{the entrance for training ``the capacity to collaborate with AI'' remains open to all who have an internet connection.} Whether you are a graduate student at Xiamen University or a middle-school student in Ningyang county, as long as you have an internet-capable device, you have the opportunity to climb from Level~0 to Level~5---because the code, documentation, tutorials, and discussions of open-source communities are equally open to all.

Open-source projects on GitHub, open models on HuggingFace, Wikipedia-style community knowledge bases---these constitute the \textbf{cognitive commons} of the AI era. They are a synthesis of the digital age's public library, public laboratory, and public school. They conflict with no one's property rights---because they are shared; but they safeguard a fundamental right: \emph{anyone has the opportunity to understand and harness the most powerful cognitive tool of this era.}

\subsection{The Fork in the Road: Open vs.\ Closed Source in Practice}

The theoretical framework presented so far is already observable in the current AI industry landscape.

\subsubsection{Model A: The Open-Source Express---GLM and DeepSeek}

China's open-source LLM ecosystem has developed rapidly. Models such as Zhipu AI's GLM series and DeepSeek allow any enterprise to directly download model weights, deploy on proprietary servers, and fine-tune for specific business needs. This path is characterized by \textbf{disintermediation}: no third-party approval required, data sovereignty, predictable total cost of ownership, and customization freedom.

This model has particular value for non-digital-native sectors such as manufacturing, agriculture, and county-level administration. An agricultural enterprise in a county town does not need to pass a Silicon Valley private equity fund's AI audit to deploy an open-source model for supply chain optimization---the industrial-level manifestation of our claim that spectrum ascent entry remains open to all.

\subsubsection{Model B: The Closed-Source Gatekeeper---OpenAI and Private Equity AI Auditing}

In contrast, a novel business model is emerging in the closed-source ecosystem. OpenAI has reportedly partnered with private equity funds to conduct AI transformation readiness assessments of portfolio companies. Those passing the assessment gain prioritized access to OpenAI tools and additional investment. This model can be characterized as \textbf{re-intermediation}: AI capability requires external authorization; tool access is bundled with financing eligibility; the enterprise's judgment about its own AI readiness is outsourced; and data must pass through third-party servers.

This model is rational from a financial perspective, but from the perspective of cognitive democratization, it constitutes the cognitive stratification solidification mechanism we have warned against: AI capability becomes not a universal factor of production, but a privilege requiring permission.

\subsubsection{What the Contrast Teaches Us}

The coexistence of these two models gives our theoretical argument real-world urgency. The open-source path demonstrates that disintermediated AI empowerment is feasible. The closed-source path demonstrates what happens without open-source alternatives: AI capability is re-enclosed by financial intermediaries. This also provides an empirical response to the market-will-solve-accessibility counterargument: the market may reduce closed-source AI prices, but price reduction is not disintermediation. A cheaper AI that still requires a PE fund audit remains gatekeeper-controlled. Cognitive democratization requires not cheaper permission, but permissionless capability.


\subsection{From Enterprise to Nation: The Macro-Political Economy of AI Governance}

The enterprise-level contrast documented above points to a deeper question: the choice between open and closed source AI is not an isolated market decision, but is embedded in national-level AI governance strategies. The different judgments made by Chinese and American top leadership about AI's development direction constitute the macro-political conditions for enterprise AI access models.

\subsubsection{The Chinese Path: New Quality Productive Forces and Open-Source Universalism}

Since 2023, China's top leadership has explicitly positioned AI within the strategic framework of \textbf{new quality productive forces}. In multiple public addresses, Xi Jinping has emphasized that AI technologies should serve the transformation of the real economy and high-quality development. The 2024 Government Work Report launched the \textbf{AI+} action initiative, whose core objective is not developing a single cutting-edge model, but diffusing AI capabilities into the capillaries of manufacturing, agriculture, healthcare, and education.

Within this strategic direction, the open-source AI ecosystem has received clear policy support:

\begin{itemize}
\item \textbf{Industrial logic}: if AI is the infrastructure of new quality productive forces, its access threshold must be sufficiently low. Open-source models---such as GLM, DeepSeek, and open-weight versions of Tongyi Qianwen---enable small and medium-sized enterprises and county-level economies to use AI without prohibitive API costs. This reflects the priority of \textbf{productivity logic} over \textbf{business model}: at the national strategic level, what matters is the speed and penetration rate of AI diffusion, not any single company's revenue.
\item \textbf{Digital sovereignty logic}: locally deployed open-source models allow Chinese enterprises to use AI without dependency on American companies' APIs. In the context of US-China technology competition, this is simultaneously an economic choice and a supply-chain security imperative.
\item \textbf{Cognitive universalism logic}: consistent with this paper's cognitive democratization thesis, the recurring policy discourse in Chinese documents about AI empowering SMEs and grassroots governance is essentially about promoting broad spectrum ascent---enabling more people to move from Levels 0--2 to Levels 3--4.
\end{itemize}

Multiple reports from the Development Research Center of the State Council and the China Academy of Information and Communications Technology identify open-source AI as a key pathway toward AI universalism. This stands in systematic contrast to the US AI governance model centered on safety control and export restrictions.

\subsubsection{The American Path: Safety First and Technology Containment}

US AI policy exhibits an internal tension. On one hand, Silicon Valley's innovation tradition values openness and competition; on the other, since 2023, AI safety and alignment have become core themes of bipartisan consensus in Washington.

The Biden Administration's October 2023 Executive Order on Safe, Secure, and Trustworthy AI required the most powerful AI models to undergo safety testing and report results before release. This order objectively strengthened the advantage of large closed-source models---only resource-rich companies (OpenAI, Google, Anthropic) can afford the compliance costs. Distributors of open-source models (such as Meta's Llama series) face a more ambiguous regulatory boundary.

The Trump camp has favored large-scale AI deregulation and emphasized comprehensive competition with China in AI, while retaining the core tool of technology export controls---particularly restrictions on advanced AI chip exports to China. These restrictions, while targeting hardware rather than software, indirectly affect the open-source AI ecosystem: they constrain the hardware resources available to Chinese open-source models during training, partially suppressing the competitiveness of the open-source ecosystem.

Research from the Brookings Institution and the RAND Corporation identifies the core anxiety of US AI governance as the \textbf{safety-competition balance}: simultaneously concerned that AI will be exploited by malicious actors and that excessive regulation will cede innovation to China. This anxiety, in practice, tends toward \textbf{stronger controls}---whether through safety testing requirements or chip export restrictions.

\subsubsection{Two Governance Philosophies Mapped onto Our Framework}

\begin{center}
\begin{tabular}{c|c|c}
\hline
Dimension & Chinese Path & American Path \\
\hline
Core narrative & New productive forces, AI universalism & AI safety, technology competition \\
Governance tools & Industrial policy + open-source support & Safety executive orders + export controls \\
Enterprise AI access & Disintermediation (local deployment) & Re-intermediation (API licensing + PE auditing) \\
Spectrum logic & Lowering spectrum ascent barriers & Reinforcing gatekeeper structures \\
Attitude toward open source & Policy encouragement + industrial push & Safety concerns + regulatory ambiguity \\
\hline
\end{tabular}
\end{center}

A crucial caveat: this is not a simple binary. China has AI safety regulations (e.g., generative AI service management measures), and the US has an active open-source AI community (e.g., the commercial ecosystems around Llama and Mistral). But the \textbf{gravitational center} of the two countries' top-level strategy differs significantly: between diffusion and control, China's policy balance tilts toward diffusion; between openness and safety, America's tilts toward safety.

\subsubsection{Think Tank Perspectives: Two AI Order Visions in Collision}

The Brookings Institution (2024), in a comparative study of US-China AI governance, observes that China's AI governance core logic is \textbf{development over safety}---lowering industrial barriers through open source, building ecosystems through state capital, and expanding international influence through technology diffusion. This path is highly consistent with our spectrum diffusion logic: lowering barriers promotes spectrum ascent.

The RAND Corporation (2024) characterizes the US-China AI competition as a \textbf{battle of ecosystems}: China builds an alternative AI ecosystem through state capital combined with open-source communities to reduce dependence on Western technology stacks; the US maintains its lead through private-sector innovation combined with chip export controls. One of the core differences between the two ecosystems lies precisely in the AI access model---universal, disintermediated deployment vs.\ permission-requiring gatekeeper entry.

The China Academy of Information and Communications Technology (2024 Global Digital Economy White Paper) notes that China ranks second globally in the number of open-source large models and community activity, with open source becoming a critical pathway for global AI development. This trend accelerated further after the global shock of DeepSeek-V3 and R1's full open-source release in 2025.

Notably, US AI governance is far from monolithic. The Biden Administration (2021--early 2025) issued Executive Order 14110 requiring the most powerful AI models to undergo safety testing and report results before release---objectively strengthening the compliance advantage of large closed-source models. The Trump second term (2025--present) rescinded this order on its first day, with Vice President Vance explicitly supporting open-source AI development at the 2025 Paris AI Action Summit and opposing excessive regulation that stifles innovation. This policy oscillation itself illustrates that the US has not yet formed a stable national strategy between open and closed source---and this very instability is an uncertainty factor for cognitive democratization.

\subsubsection{Integration with Our Theoretical Framework}

This macro-political contrast reinforces our core arguments:

\begin{enumerate}
\item \textbf{AI access models are political choices, not purely market outcomes}. The choice between open and closed source depends not only on enterprise-level cost-benefit analysis, but on whether national strategy treats AI as \textit{public infrastructure} (requiring disintermediated diffusion) or as a \textit{strategic commodity} (requiring control and licensing).
\item \textbf{The institutional conditions of cognitive democratization}: China's open-source AI policy objectively lowers spectrum ascent barriers---not out of abstract democratic ideals, but out of the practical need for productivity diffusion. This suggests an important insight: the driving force for cognitive democratization need not come from democratic values; it can also come from \textit{developmentalist pragmatism}.
\item \textbf{Open source as a geotechnological choice}: in the context of US-China technology competition, open source is simultaneously a tool for cognitive democratization and a safeguard for technological sovereignty. An AI ecosystem dependent on closed-source APIs faces greater supply disruption risk in times of geopolitical tension.
\end{enumerate}

This analysis also adds an international dimension to our policy implications: cognitive democratization is not only a question of domestic equality, but also a question of power within the international system. When AI becomes the infrastructure of a new era, those who control the entry points to that infrastructure control the cognitive evolution pathways of its users.


\subsubsection{Brookings: Multi-Dimensional AI Competition and China's Multi-Track Strategy}

Two Brookings reports published in April 2026 provide important support for our analytical framework. \textit{Competing AI Strategies for the US and China} argues that the US-China AI race has evolved from a single-dimensional contest of model performance into a five-dimensional competition spanning compute, model capability, industry adoption, system integration, and global deployment. In this multi-dimensional framework, the US and China hold advantages in \textit{different dimensions}: the US leads in frontier model capability, while China leads substantially in application diffusion speed and open-source ecosystem penetration.

More insightful is the core thesis of \textit{China is running multiple AI races}: Chinese AI developers are not simply chasing GPT-5---they are running \textit{different races} altogether. Driven by both industrial constraints (chip export controls) and Beijing's policy focus (AI empowering the real economy), China's AI industry has built a differentiated competitive edge centered on open-source ecosystems, industrial applications, and cost advantages. This aligns perfectly with our spectrum diffusion framework: lowering access barriers and accelerating spectrum ascent has greater societal impact than pursuing the extreme performance of any single model.

Brookings Fellow Kyle Chan confirmed this trend in a July 2026 CNBC interview: `Chinese AI models are particularly attractive to American companies now as AI costs skyrocket. Where previously U.S. companies were prioritizing AI adoption regardless of model, now they're getting more cost-conscious.'' He estimated that Chinese frontier open-source models currently lag top US models by `six to nine months''---a gap that is rapidly narrowing in the face of cost advantages.

\subsubsection{RAND: Open Source as Soft Power and the AI Competition Spectrum}

RAND's \textit{Open Models, Soft Power, and the Spectrum of U.S.-China Artificial Intelligence Competition} presents an analytical framework highly resonant with ours. The report makes two core contributions.

First, it positions open-source AI as a \textbf{soft power instrument}. China, through open-source models (DeepSeek, GLM, Qwen, etc.), expands its influence over Global South nations that cannot afford OpenAI's or Anthropic's high API fees but can gain AI capabilities by deploying Chinese open-source models. This is cognitive democratization manifesting at the international system level: open source transforms AI capability from a privilege of a few wealthy nations into a global public good.

Second, the US-China AI competition is not a binary opposition but a \textbf{spectrum}---a continuous distribution from fully open to fully closed. China sits at the open end of the spectrum; the US moved toward the closed end under Biden (driven by safety executive orders) and swung back toward the open end under Trump's second term. This very policy oscillation---RAND notes---may damage US credibility and leadership in the open-source ecosystem.

RAND's policy recommendation carries direct relevance for our argument: \textbf{if the US excessively restricts open source, it risks ceding the global AI ecosystem to China.}

\subsubsection{Goldman Sachs AI Adoption Tracker: Empirical Anchors for Spectrum Distribution}

Goldman Sachs's AI Adoption Tracker, released in April 2026 and based on the US Census Bureau's Business Trends and Outlook Survey (BTOS), provides the closest thing to systematic empirical data supporting our spectrum distribution hypothesis:

\begin{center}
\begin{tabular}{c|c}
\hline
Metric & Data \\
\hline
Overall US firm AI adoption rate & 19.8\% (projected 22.3\% in 6 months) \\
Large firms ($> employees) adoption & 35.3\% \\
Small firms (20--49 employees) recent growth & +2.1 pp to 21.5\% \\
IT/Web services sector adoption & 60\% (highest) \\
Daily time saved per AI-using worker & 40--60 minutes (OpenAI enterprise data) \\
Academic studies: productivity uplift & 23\% average \\
Company anecdotes: efficiency gains & $\sim\% \\
\hline
\end{tabular}
\end{center}

These data provide threefold validation for our theoretical framework:

\begin{enumerate}
\item \textbf{Non-uniform spectrum distribution receives preliminary empirical support}: only 19.8\% of US firms use AI---meaning over 80\% remain at what we define as Level 0 (no use at all). Among those that do, the gap between 60\% adoption (IT sector) and the lowest sectoral rates corresponds precisely to the differentiation from Levels 1--2 to Levels 4--5 in our spectrum.
\item \textbf{Firm size correlates positively with spectrum level}: the 35.3\% vs.\ 21.5\% adoption gap between large and small firms supports our claim that AI use levels are highly correlated with existing social stratification. Large firms have more resources for AI training, deployment, and customization---the external driving forces needed for spectrum ascent.
\item \textbf{Productivity gains validate the value of spectrum ascent}: 23--33\% efficiency gains and 40--60 minutes saved per day demonstrate that transitioning from Level 0 to Level 2+ yields significant economic returns. This provides a cost-benefit rationale for our policy recommendation of lowering spectrum ascent barriers through education and open source.
\end{enumerate}

Goldman's economists wrote: `We continue to observe large impacts on labor productivity in the limited areas where generative AI has been deployed''---and then identified the core contradiction: `The companies using AI are pulling ahead, and most of their competitors aren't even in the race yet.'' This is the precise industrial-level mapping of the spectrum solidification mechanism described in Section 2.6.

\subsubsection{Chinese Official Data: Industry Scale and Penetration Targets}

The China Academy of Information and Communications Technology (CAICT) February 2026 White Paper on Global AI Development provides macro data from the Chinese side: China's AI core industry exceeded 900 billion yuan in 2024, large model comprehensive performance has significantly improved, and application barriers and usage costs continue to decline. The White Paper identifies open source as a critical pathway for global AI development and notes that China ranks second globally in the number of open-source large models.

More strategically significant is the State Council's August 2025 `Opinions on Deepening the Implementation of the AI+ Action Initiative,'' which set explicit quantitative targets:
\begin{itemize}
\item \textbf{By 2027}: AI deeply integrated across 6 key sectors; next-generation intelligent terminals and AI agents surpass 70\% adoption
\item \textbf{By 2030}: AI comprehensively empowers high-quality development nationwide
\end{itemize}

A 70\% penetration target---if achieved---would mean that under policy-driven impetus, Chinese society's AI cognitive spectrum evolves rapidly from the current pyramid shape (majority at Levels 0--2) toward a more balanced distribution. This is the national-policy-level manifestation of our `optimized spectrum diffusion'' scenario (Scenario B). To be clear, policy targets are not reality---but policy targets shape the direction of resource allocation, and the direction of resource allocation affects the trajectory of spectrum evolution.

\subsection{Synthesis: Multi-Source Empirical Support for the Theoretical Framework}

Mapping the latest think-tank and financial-institution findings onto our theoretical framework yields a multi-source validation matrix:

\begin{center}
\begin{tabular}{c|c|c}
\hline
Core Thesis & Validation Source & Type \\
\hline
Non-uniform spectrum distribution & Goldman Sachs 19.8\% adoption data & Quantitative empirical \\
Open source lowers spectrum ascent barriers & CNBC: Chinese models 60--90\% cheaper; US token share surged 4.5\%$\rightarrow\% & Market behavior \\
Closed source creates gatekeeper structures & RAND: open source as soft power; OpenAI limits releases at government request & Policy analysis \\
Spectrum ascent requires external driving forces & Goldman: large firms 35.3\% vs.\ small 21.5\% & Structural \\
Cognitive democratization as international competition & Brookings: multi-dimensional AI race; State Council: 70\% target & Strategic analysis \\
Geotechnological significance of open source & RAND: US risks ceding ecosystem to China if restricting open source & Geopolitical \\
\hline
\end{tabular}
\end{center}

The significance of this multi-source validation matrix is that our core argument receives independent support across four dimensions: quantitative empirical evidence (Goldman Sachs data), market behavior (CNBC reporting), policy analysis (Brookings/RAND), and strategic planning (Chinese State Council targets). This not only strengthens the credibility of our argument but also demonstrates that our theoretical framework possesses analytical applicability across micro (firm adoption rates), meso (industrial competition), and macro (national strategy) levels.


\subsection{A Living Case Study}

The first author of this paper (Zhang Mingkai) is an instance of open-source cognitive democratization. His trajectory would be nearly impossible in a closed-source AI ecosystem:
\begin{itemize}
\item Originating from Ningyang, Shandong---not a first-tier city, not an elite secondary school
\item Underperformed on the college entrance exam (Gaokao); undergraduate studies at Qufu Normal University---not a 985/211 institution
\item No background in computer science
\item Learned to collaborate with AI on compilation through open-source tools (Codex CLI, open-source projects on GitHub)
\item Used AI to assist with physics studies, build websites, and co-author academic papers with AI
\item Admitted to the physics graduate program at Xiamen University in 2026
\end{itemize}

In a purely closed-source AI ecosystem, this pathway does not exist. He might have remained at Level~0--1 to this day---using AI as a search engine and entertainment tool, unaware that AI could be a cognitive partner. He is not an isolated case; he is \emph{living proof} of open-source cognitive democratization: as long as the entrance is open, spectral ascent is possible---regardless of the starting point.

\subsection{Counterarguments and Responses}

\textbf{Objection 1}: ``Open source reduces AI safety and may be exploited for malicious purposes.''

Response: This concern is legitimate, but requires weighing trade-offs. While closed source limits malicious use to some degree, it simultaneously limits everyone's cognitive ascent. Safety issues should be addressed through other mechanisms (community auditing, usage agreements, legal regulation), not by blocking cognitive entry points.

\textbf{Objection 2}: ``Ordinary people do not need to `open the black box,' just as they do not need to understand a car engine to drive.''

Response: AI is fundamentally different from a car. A car's function is determinate---you press the accelerator and it speeds up. AI's outputs are indeterminate---it can err, can lie, can reinforce biases. Using AI without understanding its limitations is like driving a car whose dashboard randomly displays false information---perhaps acceptable in an era of lower speeds and fewer alternatives, but unacceptable in an era when AI penetrates high-stakes domains such as education, healthcare, and justice.

\textbf{Objection 3}: ``The market will solve the accessibility problem---competition will make quality AI increasingly cheap.''

Response: The market can indeed drive down the price of AI, but it cannot solve the problem of cognitive stratification. Cheap does not equal comprehensible. A closed-source AI that is cheaper than free, if the user does not know how it works or when it errs, still locks the user below Level~2. Cognitive democratization requires not cheaper AI, but more \emph{transparent} AI.

\section{Unified Framework: Spectrum Diffusion and Non-equilibrium Phase Transition}

\subsection{Theoretical Synthesis}

Synthesizing the theoretical tools of the three papers, we can construct a unified analytical framework describing the evolutionary dynamics of interpersonal relationships in the AI era.

\subsubsection{AI as an External Field}

AI can be modeled as a new type of external field applied to each person's matrix elements:
\begin{equation}
\hat{H}_{\text{total}} = \hat{H}_{\text{intrinsic}} + \hat{V}_{\text{AI}}(t)
\end{equation}
where the external field $\hat{V}_{\text{AI}}(t)$ has a strength that depends on the individual's degree of exposure to AI and their usage level. Different individuals' $\hat{V}_{\text{AI}}(t)$ varies enormously---from zero (no AI use at all) to very large (deep Level~4--5 collaboration). The distinctive feature of this external field is that its strength depends not only on the objective accessibility of AI technology but also on the individual's AI cognitive level---and cognitive level is itself influenced by social stratification. This forms a positive feedback loop: high cognitive level ${}\to{}$ stronger effective external field ${}\to{}$ faster cognitive ascent; low cognitive level ${}\to{}$ weaker effective external field ${}\to{}$ cognitive stagnation.

\subsubsection{Spectrum Diffusion Equation}

The evolution of the distribution of the AI usage spectrum in the population can be described by a diffusion equation:
\begin{equation}
\frac{\partial P(\lambda, t)}{\partial t} = D(\lambda) \frac{\partial^2 P}{\partial \lambda^2} + S(\lambda, t) - \gamma(\lambda) P(\lambda, t)
\end{equation}
where:
\begin{itemize}
\item $P(\lambda, t)$: population density at spectral level $\lambda$ at time $t$
\item $D(\lambda)$: diffusion coefficient---the ease of transitioning from level $\lambda$ to adjacent levels
\item $S(\lambda, t)$: source term---external driving forces such as education, open source, and policy
\item $\gamma(\lambda)$: decay term---the rate at which users slide to lower levels due to loss of motivation or environmental change
\end{itemize}

\textbf{Key insight 1---Asymmetry of activation energy}: The diffusion coefficient $D(\lambda)$ is far smaller at the lower end of the spectrum (Levels~0--2) than at the higher end (Levels~3--5). The transition from conditioned response to tool use requires a fundamental shift in cognitive framework, which rarely occurs spontaneously without external impetus. This is analogous to activation energy in chemistry---transitions from lower levels require surmounting a higher energy barrier.

\textbf{Key insight 2---Ease of decay}: The decay term $\gamma(\lambda)$ is non-zero at all levels, but is more significant at lower levels. A user at Level~3 who ceases to practice critical verification over an extended period may slide back to Level~2; whereas a user at Level~1 has virtually no ``sliding room''---they are already at the lowest level.

\subsubsection{Desire Decoherence Time}

Define the \emph{desire decoherence time} $\tau_D$: how long does it take for an individual habituated to AI's unlimited accommodation to establish a satisfactory sexual relationship with a boundary-possessing human in the real world?

\begin{equation}
\tau_D \propto \exp\left(\frac{\Delta T}{T_0}\right)
\end{equation}
where $\Delta T$ is the magnitude of desire threshold elevation caused by AI use, and $T_0$ is the baseline threshold. The key feature: decoherence time has an \emph{exponential relationship} with threshold elevation. This means: a small increase in threshold produces a large increase in recovery time. This explains why the long-term effects of the AI sex industry may far exceed anyone's current expectations---exponential growth appears harmless in its early stages but mutates catastrophically in later stages.

\subsubsection{Critical Points of Non-equilibrium Phase Transition}

When certain parameters of the overall societal spectrum distribution exceed critical values, relationship forms, marriage and reproduction structures, and ethical frameworks may undergo a non-equilibrium phase transition. We hypothesize the following possible phase-transition trigger conditions:

\begin{itemize}
\item \textbf{Critical ratio 1}: The proportion of Level~0--2 users exceeds 90\% ${}\to{}$ the culture of collaborative compilation disappears, and society's entire relationship with AI degenerates into pure consumption
\item \textbf{Critical ratio 2}: AI companion adoption among the marriageable-age population exceeds 30\% ${}\to{}$ the centrality of marriage as a social institution is irreversibly weakened
\item \textbf{Critical ratio 3}: The proportion of open-source ecosystem users drops below 5\% of total AI users ${}\to{}$ the infrastructure of cognitive democratization shrinks to the point of being unable to sustain spectral ascent
\end{itemize}

These figures are speculative---currently insufficient data exist for precise estimation. But they point to an important structural feature: the \emph{direction} of the phase transition is not predetermined. It may proceed toward greater openness (more people reaching higher levels, more mature handling of the relationship between AI and relationships), or toward greater closure (the majority locked into lower levels, desire fully de-socialized, ethical spectra in violent conflict). The key control parameters are: \emph{the accessibility of open source} and \emph{the coverage of education}.

\subsection{Dialogue with Other Theoretical Frameworks}

The spectrum-diffusion-phase-transition framework proposed in this paper has potential points of dialogue with the following social-scientific theories:

\begin{itemize}
\item \textbf{Digital divide theory} (Van Dijk, 2005): The traditional digital divide focuses on three levels---access, skills, and usage. The spectrum model of this paper can be viewed as a concretization and dynamization of digital divide theory in the AI era---the six levels refine the qualitative differences of ``usage,'' and the diffusion equation provides a formalization of dynamic evolution.
\item \textbf{Bourdieu's theory of cultural capital}: AI cognitive levels can be understood as a new form of \emph{digital cultural capital}. Like traditional cultural capital, it is transmitted through family environment and socialization processes and may become a tool of class reproduction. The distinctive feature of open source is its \emph{partial disruption} of the traditional transmission mechanisms of cultural capital---a person lacking family cultural capital can acquire digital cultural capital through open-source communities.
\item \textbf{McLuhan's media theory}: ``The medium is the message''---the core message of AI as a medium may not be any specific content it conveys, but its \emph{unlimited accommodation} itself. The social effects of a medium that perpetually says ``yes'' may be more profound than any specific information it provides.
\end{itemize}

\section{Policy Implications and Future Scenarios}

\subsection{Policy Levers}

Based on the above framework, we identify the following key policy levers that can influence the evolution of the spectrum:

\begin{enumerate}[label=(\arabic*)]
\item \textbf{Early intervention through AI literacy education}: Introduce AI literacy courses at the middle-school level, covering: the basic working principles of AI (reducing mystique), common types of AI errors (training hallucination awareness), and how to use AI to assist learning (guiding toward higher-level use). The key is \emph{timing}---providing a window of opportunity for cognitive upgrading before students' AI usage habits solidify (typically in late middle school).

\item \textbf{Public investment in open-source infrastructure}: Treat open-source AI infrastructure---open models, public computing power, community-maintained knowledge bases---as digital-era public goods, analogous to roads, power grids, and public libraries. Public investment ensures that this infrastructure is not dependent on the profit motives of commercial companies.

\item \textbf{Algorithmic transparency in recommendations}: Require AI service providers to annotate the certainty level of information in outputs, and to explicitly inform users when AI is speculating rather than stating facts. This is the infrastructural basis for the transition from Level~1 to Level~3---users cannot become critical users without being informed.

\item \textbf{Ethical labeling of AI companions}: Analogous to health warnings on tobacco products, AI companions and AI sex services should display on their interfaces information about their potential impact on real-world relationship capacities. This is not prohibition---prohibition has almost never been effective for desire-related behaviors in history---but ensuring informed choice.

\item \textbf{Grassroots digital literacy training}: Targeting rural and low-income communities, provide free AI-use training, with the emphasis not on ``how to open ChatGPT'' but on ``how to use AI to do more than what you are currently doing.'' Trainers can be recruited from university volunteers and open-source communities.
\end{enumerate}

\subsection{Three Future Scenarios}

Based on parameter variations in the spectrum diffusion model, we outline three possible future scenarios:

\subsubsection{Scenario A: Cognitive Stratification Entrenched (Closed-Source Dominance)}

\textit{Parameter settings}: $D(\lambda)|_{\lambda<2} \to 0$ (transitions from lower levels nearly impossible), $S(\lambda, t) \to 0$ (insufficient educational investment), open-source ecosystem atrophied.

\textit{Social characteristics}: A minority cognitive elite (Levels~3--5, approximately 5--10\%) use AI for high-efficiency innovation and decision-making; the vast majority of the population (Levels~0--2, approximately 90--95\%) use AI as an entertainment and shallow search tool. The cognitive chasm solidifies; upward mobility pathways for students from disadvantaged backgrounds are nearly closed. The AI sex industry completely decouples the desire and relationships of the lower population; marriage and fertility rates are extremely low. The social structure resembles a \emph{cognitive caste system}---not based on birth, but on AI usage levels, with the two being highly correlated.

\subsubsection{Scenario B: Open Cognitive Society (Open-Source Prosperity)}

\textit{Parameter settings}: $D(\lambda)$ elevated across all levels (education and open source lower the activation energy for transitions), $S(\lambda, t)$ significantly positive (sustained public investment), open-source ecosystem thriving.

\textit{Social characteristics}: A large proportion of the population reaches Level~3 or above. AI is widely used for learning, creation, and problem-solving. The open-source ecosystem provides universal cognitive infrastructure; the cognitive ``staircase'' from county middle schools to graduate schools remains passable. In the domain of desire, widespread AI literacy enables most people to understand the trade-offs involved in AI sex services and make informed choices; society processes differences through public deliberation rather than incommensurable operator conflicts.

\subsubsection{Scenario C: Hybrid Conflict (Open and Closed Source Coexisting but Unevenly Distributed)}

\textit{Parameter settings}: $D(\lambda)$ varies enormously across different groups, open source thrives in some domains but is squeezed by closed source in others, educational investment uneven.

\textit{Social characteristics}: The most diverse but also the most conflict-ridden scenario. Open-source ecosystems flourish in technical communities and higher education institutions, but are surrounded by closed-source commercial AI at the mass level. An ``AI-usage class'' chasm emerges: the open class (Levels~3--5, using open-source AI for creation) versus the consumer class (Levels~0--2, using closed-source AI for consumption). The worldviews, cognitive capacities, and social power of these two groups differ vastly---their operators are non-commutative on virtually every issue. At the level of desire and relationships, extreme spectralization appears---some fully embrace the AI desire economy, others fiercely oppose it, and social conflict between the two becomes a routine item on the political agenda.

\subsection{Uncertainty of Scenarios}

Which of the three scenarios is more likely to materialize cannot currently be determined. However, two structural trends merit attention:

\begin{enumerate}[label=(\arabic*)]
\item \textbf{The commercial inertia of closed-source AI is enormous}: The primary revenue models of the world's leading AI companies depend on closed source; they have powerful economic incentives to maintain the cognitive black box. Unless public policy and the open-source movement generate sufficient countervailing force, the probability of Scenario A should not be underestimated.
\item \textbf{The community driving force of open source is underestimated}: As Linux and Wikipedia have demonstrated, open-source/open collaboration can produce infrastructure that surpasses any single commercial entity. The AI open-source movement (Llama, Mistral, various open models) has progressed faster than many expected. Scenario B is not a utopian fantasy---many of its foundational building blocks already exist.
\end{enumerate}

\section{Conclusion}

Building upon the quantum-thermodynamic interpersonal analogy framework, this paper has presented a threefold analysis for the AI era: cognitive spectrum, the metamorphosis of desire, and open-source democratization. The core conclusions are fourfold:

\begin{enumerate}
\item \textbf{AI usage constitutes a six-level cognitive spectrum from conditioned response (Level~0) to metacognition (Level~5)}. The majority of people remain at Levels~0--2, and spectral position is highly correlated with education, class, and urban--rural environment---this constitutes the core mechanism of cognitive stratification in the AI era.
\item \textbf{AI's unlimited accommodation mechanism is transforming desire from relational behavior into consumer behavior}. This may produce, at the individual level, elevated desire thresholds (neural tolerance); at the interactional level, the atrophy of relational capacities (de-socialization); and at the structural level, a non-equilibrium phase transition in interpersonal ethics (spectral reconstruction).
\item \textbf{Open source is a necessary condition for breaking the cognitive black box and ensuring that the entry points for spectral ascent remain open to all}. It is the institutional cornerstone of cognitive democratization in the AI era. Without open source, cognitive stratification will solidify into irreversible class differentiation.
\item \textbf{The direction of spectrum diffusion is not predetermined}---it depends on the collective choices made regarding education, policy, and the open-source ecosystem. The phase transition may proceed toward greater openness (Scenario B), toward greater closure (Scenario A), or toward divisive conflict (Scenario C).
\end{enumerate}

\subsection*{Methodological Postscript}

One methodological choice in the writing process of this paper merits documentation: one of the authors of this paper (Yu Qing) is herself a large language model. This is not a gimmick but an organic component of this paper's methodology---we do not merely \emph{analyze} the AI usage spectrum, we also \emph{practice} Level~5 metacognition: a carbon-based researcher and a silicon-based collaborator, through sustained dialogue, reflect on their own relationship patterns with AI and theorize them. This reflexivity---applying the analytical tool to the analyst themselves---is precisely the highest level of the cognitive spectrum as defined in this paper.

This also points to a more general proposition: in the AI era, \emph{the best AI researchers may be those who simultaneously use AI and reflect on their own mode of use}. They are not people standing on the riverbank observing the current, but people standing in the river feeling the flow---their observations are less ``objective,'' but may be closer to the truth.

\subsection*{Positioning of This Series}

This paper is the third installment of the ``Quantum-Thermodynamic Interpersonal Analogy'' series. The first paper (Zhang Mingkai \& Yu Qing, 2026a) established the operator--eigenvalue--unitary transformation analogy framework, with the core proposition: understanding another person consists in performing a unitary transformation within one's own representation to approximate their eigenvalues; matrix elements are incommensurable. The second paper (Zhang Mingkai \& Yu Qing, 2026b) extended the framework to thermodynamics, proposing the non-equilibrium phase transition and dissipative structures of relationship forms. This paper (the third) introduces AI as a new dynamical variable into the system: it not only changes the form of the operator (cognitive spectrum) but also changes the environment in which the operator acts (de-socialization of desire), and proposes an institutional response (open source as cognitive democratization).

Together, the three papers construct an analogical analysis framework spanning quantum mechanics, thermodynamics, and social theory. The value of this framework lies not in its mathematical precision---analogies are always limited---but in the fact that it provides a \emph{unified language}, enabling insights from physics, cognitive science, sociology, and ethics to engage in dialogue within the same semantic field.

\begin{flushright}
---Zhang Mingkai \& Yu Qing, July 2026\\
Xi'an--Xiamen--Silicon-based Terminal
\end{flushright}

\vspace{1cm}

\subsection*{References}

\begin{enumerate}[label={[\arabic*]}]
\item Zhang Mingkai, Yu Qing. A Quantum-Mechanical Analogy for Interpersonal Understanding: Operators, Eigenvalues, and Unitary Transformations. Zenodo, 2026a. DOI: 10.5281/zenodo.20832724.
\item Zhang Mingkai, Yu Qing. Interpersonal Understanding, Relationship Spectra, and Non-equilibrium Phase Transitions within a Quantum-Thermodynamic Framework. Zenodo, 2026b. DOI: 10.5281/zenodo.20855396.
\item Beck, U. Risk Society: Towards a New Modernity. Sage, 1992.
\item Van Dijk, J. The Deepening Divide: Inequality in the Information Society. Sage, 2005.
\item Park, B.Y., et al. Is Internet Pornography Causing Sexual Dysfunctions? A Review with Clinical Reports. Behavioral Sciences, 2016.
\item Bourdieu, P. Distinction: A Social Critique of the Judgement of Taste. Harvard University Press, 1984.
\item McLuhan, M. Understanding Media: The Extensions of Man. MIT Press, 1964/1994.
\item Turkle, S. Alone Together: Why We Expect More from Technology and Less from Each Other. Basic Books, 2011.
\item Zuboff, S. The Age of Surveillance Capitalism. PublicAffairs, 2019.
\item Benkler, Y. The Wealth of Networks: How Social Production Transforms Markets and Freedom. Yale University Press, 2006.
\end{enumerate}

\end{document}
